\section{Introduction}
\label{sec:intro}

The explosive growth in specialized, pre-trained neural network models has catalyzed a paradigm shift in machine learning engineering: rather than training a monolithic, multi-task network from scratch, researchers and practitioners increasingly opt to fuse specialized ``expert'' models. This practice, known as \emph{model merging}, combines the weights of multiple fine-tuned models to achieve multi-task capability with zero additional training cost and no expansion of the physical deployment footprint \cite{Wortsman2022}. Recently, \emph{dynamic model merging} (or test-time expert blending) has emerged as an active frontier. Unlike static merging, which computes a single fixed set of interpolation coefficients, dynamic model merging utilizes an input-dependent router to blend expert weights (such as LoRA adapters) on the fly at inference time.

Despite its theoretical promise, deploying dynamic routers in real-world production streams reveals severe empirical vulnerabilities that are often overlooked in idealized laboratory settings. As empirically driven researchers, we believe that true progress in machine learning comes from rigorous stress-testing and exhaustive empirical validation across realistic deployment conditions. In this work, we zoom in on two critical, interconnected failure modes that compromise the stability and reliability of dynamic model merging systems:

\begin{enumerate}
    \item \textbf{Calibration Data Scarcity (The Small-$N$ Overfitting Regime):} Standard parametric routers rely on a small set of labeled calibration samples to learn task-routing coefficients via gradient descent. When this calibration set is scarce ($N \le 32$ samples), parametric models overfit severely. The resulting transductive collapse leads to high variance across training seeds and terrible generalization on test data.
    \item \textbf{Deployment Stream Batch Heterogeneity (Heterogeneity Collapse):} Standard dynamic routers assume homogeneous batches at inference time. However, real-world deployment streams are highly heterogeneous, containing mixed tasks within the same execution batch. Processing these mixed batches causes the representations of different tasks to average out. The router, lacking sample-wise isolation, outputs flat, near-uniform routing weights across experts, a catastrophic phenomenon we refer to as \emph{heterogeneity collapse}.
\end{enumerate}

To overcome these vulnerabilities, we analyze a robust, dual-pathway routing design pattern called \textbf{Confidence-Gated Hybrid Routing (CGHR)}, augmented by a stream-partitioning technique named \textbf{Micro-Batch Homogenization (MBH)}. 
CGHR introduces a confidence-driven gating mechanism that continuously balances the expressive flexibility of a trained parametric linear router (Pathway A) against the zero-shot stability of a completely non-parametric, parameter-free subspace router (Pathway B). Specifically, when the parametric router's prediction confidence is high, CGHR leverages its learned task-routing weights. Conversely, when confidence falls below a threshold $\gamma_{\text{conf}}$, indicating an out-of-distribution (OOD) task or an ambiguous, scarce calibration sample, the system automatically falls back to the parameter-free subspace routing (PFSR) pathway, which projects the penultimate features onto expert classification manifolds and requires zero training data.

To resolve heterogeneity collapse, MBH dynamically partitions incoming mixed-task batches into homogeneous micro-batches on the fly, calculating routing coefficients independently and executing localized weight fusion.

Rather than relying on minor theoretical justifications, we ground this work in systematic, simulation-based evaluations. Utilizing a synthetic multi-task simulation sandbox pre-calibrated to recreate realistic expert ceilings across MNIST, Fashion-MNIST, CIFAR-10, and SVHN, we execute hundreds of parallel parameter sweeps over multiple seeds and deployment configurations. Our sweeps over the confidence threshold $\gamma_{\text{conf}}$, calibration sample complexity $N$, and batch size $B$ demonstrate that:
\begin{itemize}
    \item CGHR discovers a robust \emph{peak performance envelope} at intermediate thresholds ($\gamma_{\text{conf}} \approx 0.85$ using Maximum Probability), outperforming both pure parametric and non-parametric baselines.
    \item Under severe calibration data scarcity ($N = 16$), CGHR maintains near-perfect, flatline-like stability, leveraging the non-parametric fallback to completely bypass transductive overfitting.
    \item The integration of MBH completely shields the system from heterogeneity collapse, delivering perfectly flat performance curves across all batch sizes (up to $B=512$).
\end{itemize}

While the Isolating Coordinate Sandbox is a simplified, 1-layer coordinate simulation, it acts as a highly controlled mathematical instrument to isolate and audit individual routing behaviors. We discuss its structural limitations and systems-level overheads in detail to provide an honest, deployment-aware perspective.

The rest of the paper is organized as follows: Section \ref{sec:related} reviews related work and identifies key empirical gaps. Section \ref{sec:method} presents the mathematical formulation of CGHR and MBH. Section \ref{sec:experiments} details our exhaustive experimental results, sensitivity sweeps, and limitations. Section \ref{sec:conclusion} concludes the paper and outlines future directions for real-world scaling.
